\clearpage
\section{Copying robustness for every original closing context}
\label{sec:copying-contextual-robustness}
The concrete copying robustness theorem now concludes contextual refinement
for the entire original observation class. A caller no longer needs to supply
a recursive filling certificate for each observation. We construct that
certificate from the context's existing sole-hole address, prove that its
results are recursively isomorphic to the actual addressed fillings, and
identify their initial states in the proved representation-class machine.
This closes the addressed-closing-context connection left open in
Section~\ref{sec:copying-specimens}.

The conclusion uses \texttt{Combined.execution.Refines} and
\texttt{Combined.execution.Equiv} on actual open program graphs, for arbitrary
variable and atom interfaces. It quantifies over all original binding-free
closing contexts, including arbitrary surrounding depth and nested box holes.
Forward and converse robustness retain their quantitative cost relations.
Only at the termination corollary do we specialize to the universal cost
relation. The original Proposition~12.2 terms, Figure~16 translation and all
machine rules are unchanged.

This is the concrete observation-level use of PDF Theorem~8.11 (HTML
Theorem~6, \href{https://arxiv.org/html/1907.01257v5\#S8.SS4}{Section~8.4}).
It retains output closure, input safety and robustness as hypotheses on a
template. It does not prove any of the eleven template families automatically.
Correspondence with explicitly focused paper graphs, unique-hole inventories
and replacement tuples remains separate. In particular, the current specimen
interface ranges over binding-free contexts with recursive certificates;
relating that domain to the paper's tuple-based template obligations is still
required. No final example equivalence is asserted here.

\Needspace{10\baselineskip}
\subsection{Collapsing copies of retained edges}
\termdef{copying-corolla-collapse}
A \emph{corolla collapse} removes the auxiliary carrier names introduced by
expanding an edge into a literal copy of itself. A \texttt{LabelCopy}
certificate says that there is exactly one edge and that its source and
target lists are the full ordered boundaries. Such a component has no internal
vertices: source coverage would place an internal vertex on the sole edge's
source list, whereas every entry on that list is a boundary input. These are
disjoint constructors of the vertex type.

\lean{Copying/FillingCollapse}{9}{14}{LabelCopy.no_internal}

For a host $C$ and a family of certified corollas $F_e$, an internal vertex of
$C[F]$ is either an original host vertex or an internal vertex of some $F_e$.
The latter case is impossible. Thus the vertex bijection keeps the host
vertex and its inverse inserts the host summand. An edge is a pair $(e,x)$;
its image is $e$, and the inverse chooses the certified sole edge of $F_e$.
The certificate's uniqueness equation proves that these functions are inverse.

\lean{Copying/FillingCollapse}{22}{26}{FillingCollapse.vertices}
\lean{Copying/FillingCollapse}{28}{32}{FillingCollapse.edges}

To check ordered incidence, substitute the component's full boundary list
into the expansion source or target formula. Each boundary input is attached
to the corresponding original source port, and each boundary output to the
corresponding original target port. The collapse cancels the host embedding.
Mapping the finite list of all boundary indices therefore returns the
original list, in its original order. This proves both equations, including
nullary edges and repeated wire types.

\lean{Copying/FillingCollapse}{34}{42}{FillingCollapse.source}
\lean{Copying/FillingCollapse}{44}{52}{FillingCollapse.target}

\Needspace{8\baselineskip}
If every copied label is the original label, these bijections and incidence
equations give an ordinary graph isomorphism $C[F]\cong C$. More generally,
if each copied label is a prescribed replacement label of the same ordered
signature, they give $C[F]\cong C[\ell]$, where $C[\ell]$ is the host with
those edge labels replaced. The edge-label equation follows from the sole-edge
certificate. No semantic equivalence or box permutation is used.

\lean{Copying/FillingCollapse}{55}{61}{FillingCollapse.identity}
\lean{Copying/FillingCollapse}{64}{72}{FillingCollapse.relabel}

\Needspace{12\baselineskip}
\subsection{One addressed replacement and its enclosing boxes}
\termdef{copying-single-expansion}
The \emph{single-replacement expansion family} at edge $e$ uses the actual
replacement $H$ at $e$ and a literal label corolla at every other edge.
The dependent equality in the selected branch preserves its exact signature.
At $e$ the family is $H$; away from $e$ it has a retained-copy certificate.

\lean{Copying/FillingCollapse}{80}{83}{SingleExpansion.family}
\lean{Copying/FillingCollapse}{85}{87}{SingleExpansion.selected}
\lean{Copying/FillingCollapse}{89}{91}{SingleExpansion.retained}

The simultaneous expansion of this family is recursively isomorphic to the
original single-edge plug $C[e:=H]$. To prove this, peel the selected expansion
out using the generic expansion-peeling isomorphism. The residual expansion
copies all surviving host edges and all inserted edges. Corolla collapse
removes precisely that residual expansion. Composing the two isomorphisms
proves the claim, with the original inserted graph intact.

\lean{Copying/FillingCollapse}{94}{103}{SingleExpansion.iso}

At a box address, the operation changes the box body while retaining its
binding count and variable and atom interface. A literal box containing the
new body has the same ordered signature as the old box. We first transport
its boundary equality, then use it as the selected expansion component.
Boundary transport preserves a label-copy certificate: eliminate the two
interface equalities and use the identity-boundary equation.

\lean{Copying/AddressedFillings}{8}{13}{LabelCopy.reboundary}
\lean{Copying/AddressedFillings}{20}{22}{BoxExpansion.component}
\lean{Copying/AddressedFillings}{24}{25}{BoxExpansion.copy}
\lean{Copying/AddressedFillings}{27}{27}{BoxExpansion.family}

Corolla collapse with the prescribed box label gives a recursive isomorphism
to the original \texttt{replaceBox} construction.
All other labels and all outer incidence lists are retained. If the two new
bodies already have a recursive paired-filling certificate, the selected edge
uses its boxed constructor; every other edge uses the retained constructor.
Thus paired filling propagates through an enclosing original box.

\lean{Copying/AddressedFillings}{31}{50}{BoxExpansion.iso}
\lean{Copying/AddressedFillings}{53}{62}{BoxExpansion.paired}

\termdef{copying-addressed-certificate}
An \emph{addressed filling certificate} for a template pair $(G,H)$ and an
original sole-hole site $s$ consists of graphs $L,R$, a recursive paired
filling of the original context producing $L,R$, and isomorphisms
\[
 L\cong \operatorname{plugAt}(s,G),\qquad
 R\cong \operatorname{plugAt}(s,H).
\]
We prove its existence by induction on the site. At a shallow hole, use the
single-replacement families and the template pair at the selected edge.
At an address inside a box, apply the induction hypothesis to its body,
propagate the resulting paired certificate through the box, and compose
box-expansion collapse with recursive isomorphism under body replacement.
This induction follows the actual address; it never substitutes a different
observation context.

\lean{Copying/AddressedFillings}{69}{90}{AddressedFillings.paired}

For a closing context and hole-free open programs, the existing sole-hole
filling theorem makes both actual results hole-free. Recursive isomorphism
reflects this property to $L,R$. Consequently the certificate has all the
global hole-freedom proofs needed by the concrete specimen interface.
Focus freedom follows from the pre-template and the original context.

\lean{Copying/AddressedFillings}{94}{100}{AddressedFillings.closing}

\Needspace{12\baselineskip}
\subsection{Initial lifting and contextual robustness}
\termdef{copying-closing-lifting}
The \emph{closing-context initial lifting} is a rooted, remote specimen witness
whose projections are the actual initial observations of $C[G]$ and $C[H]$.
First, any recursive isomorphism of whole closed graphs maps their search
input-zero coordinates to one another. It therefore yields a state
representation relation and equality of their representation classes. This
is an equality of machine states, not an added transition or a cost allowance.

\lean{Copying/ContextualRobustness}{10}{12}{StateClass.initial_iso}

The allowed-context predicate is exactly binding freedom. For every original
closing context $C$, extract the preceding addressed certificate, replace
its actual initial states by the equal class representatives, and expose the
outer expansion constructor. The context's binding-freedom and focus-freedom
proofs, paired components and hole-free results construct a concrete specimen.
The previously proved initial-lifting theorem supplies its rooted remote
witness. Transport back along the class equalities proves the desired lifting
for the original observations.

\lean{Copying/ContextualRobustness}{17}{19}{Specimens.bindingFreeContexts}
\lean{Copying/ContextualRobustness}{24}{37}{Specimens.closing_initial_lifting}

\termdef{copying-contextual-robustness}
The \emph{concrete contextual robustness theorem} now says the following.
Let $T$ be output closed, let $(Q,P,R)$ be a reasonable cost triple, and
suppose the concrete specimen system for $T$ is input safe and robust for
these costs. If the actual open program graphs $G,H$ are a pair of $T$, then
$G$ contextually refines $H$ for $Q$. For each original closing context,
apply forward sufficiency to its constructed initial lifting. There is no
remaining premise about determinism, refocusing, padding invariance, locality,
or existence of observation certificates.

\lean{Copying/ContextualRobustness}{41}{46}{Specimens.contextual_refines}

The converse uses the opposite specimen system and its robustness proof.
Its input-safety condition uses $Q^{-1}$, exactly as in the abstract theorem;
it is not replaced by forward safety for arbitrary costs. Combining these
two conclusions gives quantitative contextual equivalence of the actual open
graphs. Each direction ranges over the complete original observation class.

\Needspace{4\baselineskip}
\lean{Copying/ContextualRobustness}{49}{55}{Specimens.contextual_refines_converse}
\lean{Copying/ContextualRobustness}{59}{67}{Specimens.contextual_equiv}

Specialize to the proved reasonable triple consisting of universal final
cost, improvement on the left, and universal right auxiliary cost. The
universal relation equals its converse. The observation characterization then
gives termination of $C[G]$ if and only if termination of $C[H]$ for every
original $C$. No assumption that callbacks or either computation terminate is
introduced. Stuckness remains a negative observation, as before.

\lean{Copying/ContextualRobustness}{71}{78}{Specimens.contextual_termination}

For terms, require that their original translated graphs belong to the
proved template at each adequate common depth. Apply the contextual theorem
at each such depth. This gives exactly the existing term observational
equivalence, with both variable and atom environments unrestricted.
The template membership hypothesis is an application obligation, not a
translation-correctness axiom or the desired observational equivalence itself.

\lean{Copying/ContextualRobustness}{82}{91}{Specimens.term_equiv}

\Needspace{12\baselineskip}
\subsection{Translated contexts, cofinal allowances and regressions}
For any original closing term context, the previously proved graph-filling
law supplies a graph context with the actual translated host, a sole hole and
binding freedom. Packaging it as a closing context retains literal host
identity and the reconstruction law: supplying the suitably padded original
fragment returns the interpretation of the original filled term, up to
recursive isomorphism at equal depth. Bound variables remain in the annotated
hole's interface. Combining this law with closing-context initial lifting
gives actual specimen witnesses for all such translated contexts.

\lean{Copying/TranslatedSpecimens}{11}{18}{ContextTranslation.closing_context}
\lean{Copying/TranslatedSpecimens}{23}{33}{ContextTranslation.closing_specimens}

One can weaken the term membership premise to arbitrarily large common
allowances. Given any lower bound, choose the related translated pair above
that bound and apply quantitative contextual robustness there. The concrete
cofinal-observation theorem then gives equivalence at every adequate allowance,
using the already proved exact padding reflection. Thus the argument imposes
no new depth-stability axiom on the template.

\lean{Copying/TranslatedSpecimens}{40}{52}{Specimens.term_equiv_cofinal}

Finally, every original term context preserves the resulting equivalence.
First derive equivalence of the two terms by the template theorem, then apply
the proved concrete term-context congruence. This works beneath binders and
when the terms have different minimal translation depths. Both isomorphism
and padding premises of the older abstract congruence result are already
discharged by the concrete copying execution.

\lean{Copying/TranslatedSpecimens}{57}{65}{Specimens.term_context_equiv}

The first regression constructs initial lifting for arbitrary numeral pairs
in every original closing context, at every allowance. This deliberately
checks the structural bridge only; it does not assert robustness or
equivalence of arbitrary numerals.

\lean{Copying/ContextualRobustnessExamples}{11}{16}{ContextualRobustnessExamples.numbers_initial}

A second deliberately broad pre-template relates any two focus-free graphs.
It is suitable for testing structural certificates but is not output closed:
relate the negation corolla to itself and apply the already checked positive
operation path from its input to its output. Neither side is one way.
This negative regression prevents treating structural certificate construction
as a proof of the semantic template premises.

\lean{Copying/ContextualRobustnessExamples}{19}{21}{ContextualRobustnessExamples.focusFreePairs}
\lean{Copying/ContextualRobustnessExamples}{24}{29}{ContextualRobustnessExamples.focusFreePairs_not_closed}

The nested-binder regression uses the actual context $\lambda.\lambda.[-]$,
whose hole has two bound variables in scope. It fills that hole with index
zero on the left and index one on the right. The addressed bridge gives a
paired certificate, and composing each side's isomorphism with the original
translation-filling law recovers its respective unchanged term translation.
The test does not identify these two terms or claim they are equivalent.
A final check derives the sourced-hole property for every closing context
from its sole-hole address and its nonempty term-input interface. The
zero-input exception to Lemma~E.1 has not been erased.

\lean{Copying/ContextualRobustnessExamples}{34}{54}{ContextualRobustnessExamples.nested_binders}
\lean{Copying/ContextualRobustnessExamples}{57}{59}{ContextualRobustnessExamples.sourced_hole_required}

Appendix~B.1 now omits the completed closing-context certificate and initial
lifting obligation. Four broader work packages remain, including the other
paper-correspondence checks, the eleven template families, semantic weakening
and Figure~22. The estimate remains 60--70\% of total mathematical effort,
with low confidence. This milestone removes a universal observation obligation;
it does not justify estimating the remaining template proofs by declaration
counts.
